\section{Conclusion de l'état de l'art}

En conclusion, cet état de l’art met en évidence des travaux proches du projet eSWARM ainsi que les concepts clés du contrôle d’essaims et de la commande événementielle. Le système ModQuad~\cite{gabrich_modquad-dof_2020} apparaît comme le plus proche, mais il repose sur une logique partiellement centralisée, à l’opposé de l’approche distribuée visée pour eSWARM. À l’inverse, les travaux du laboratoire de Tokyo~\cite{sugihara_beatle_2024,sugihara_design_2023} sont décentralisés, mais s’appuient sur des drones nettement plus sophistiqués que ceux actuellement employés dans eSWARM.
Par ailleurs, la commande événementielle constitue un domaine de recherche actif, riche en approches. Toutefois, les applications concrètes restent rares dans le contexte des essaims de robots modulaires et reconfigurables, ce qui rend le projet eSWARM unique.